% =============================================================================
% 图 2：异步 All-to-Allv 流水时序
% 编译：xelatex fig2_async_pipeline.tex
% =============================================================================
\documentclass[tikz,border=4pt]{standalone}
% =============================================================================
% DCC-KV 论文图表 —— 共享前言
%
% 由各 fig*.tex 通过 % =============================================================================
% DCC-KV 论文图表 —— 共享前言
%
% 由各 fig*.tex 通过 % =============================================================================
% DCC-KV 论文图表 —— 共享前言
%
% 由各 fig*.tex 通过 \input{preamble} 载入，必须位于 \documentclass 之后。
%
% 编译：在 figures/ 目录下执行 xelatex fig1_architecture.tex
%      （含中文，必须用 xelatex；也可在 paper/ 下执行 make figures）
%
% 字体：默认使用 Windows 简体中文环境自带字体。
%      在 Linux/macOS 上请把下面三行替换为该系统已安装的 CJK 字体，例如
%      \setCJKmainfont{Noto Serif CJK SC}[BoldFont=Noto Sans CJK SC]
% =============================================================================

\usepackage{xeCJK}
\setCJKmainfont{SimSun}[BoldFont=SimHei]
\setCJKsansfont{Microsoft YaHei}

\usepackage{tikz}
\usepackage{amsmath,amssymb}
\usetikzlibrary{arrows.meta,positioning,fit,backgrounds,calc,shapes.geometric,decorations.pathreplacing}

% --- 统一样式 ---------------------------------------------------------------
\tikzset{
  dev/.style    = {draw=black!70, fill=blue!6,  rounded corners=2pt, minimum width=26mm,
                   minimum height=11mm, align=center, font=\small, line width=0.5pt},
  devr/.style   = {draw=black!70, fill=orange!10, rounded corners=2pt, minimum width=26mm,
                   minimum height=11mm, align=center, font=\small, line width=0.5pt},
  blk/.style    = {draw=black!60, fill=gray!8, rounded corners=1.5pt, minimum width=17mm,
                   minimum height=8mm, align=center, font=\scriptsize, line width=0.4pt},
  msg/.style    = {draw=black!60, fill=green!8, rounded corners=1.5pt, minimum width=15mm,
                   minimum height=7mm, align=center, font=\scriptsize, line width=0.4pt},
  flow/.style   = {-{Stealth[length=2mm]}, line width=0.7pt, draw=black!75},
  flowd/.style  = {-{Stealth[length=2mm]}, line width=0.7pt, draw=black!50, dashed},
  lab/.style    = {font=\scriptsize, inner sep=1.5pt, fill=white, fill opacity=0.85,
                   text opacity=1},
}
 载入，必须位于 \documentclass 之后。
%
% 编译：在 figures/ 目录下执行 xelatex fig1_architecture.tex
%      （含中文，必须用 xelatex；也可在 paper/ 下执行 make figures）
%
% 字体：默认使用 Windows 简体中文环境自带字体。
%      在 Linux/macOS 上请把下面三行替换为该系统已安装的 CJK 字体，例如
%      \setCJKmainfont{Noto Serif CJK SC}[BoldFont=Noto Sans CJK SC]
% =============================================================================

\usepackage{xeCJK}
\setCJKmainfont{SimSun}[BoldFont=SimHei]
\setCJKsansfont{Microsoft YaHei}

\usepackage{tikz}
\usepackage{amsmath,amssymb}
\usetikzlibrary{arrows.meta,positioning,fit,backgrounds,calc,shapes.geometric,decorations.pathreplacing}

% --- 统一样式 ---------------------------------------------------------------
\tikzset{
  dev/.style    = {draw=black!70, fill=blue!6,  rounded corners=2pt, minimum width=26mm,
                   minimum height=11mm, align=center, font=\small, line width=0.5pt},
  devr/.style   = {draw=black!70, fill=orange!10, rounded corners=2pt, minimum width=26mm,
                   minimum height=11mm, align=center, font=\small, line width=0.5pt},
  blk/.style    = {draw=black!60, fill=gray!8, rounded corners=1.5pt, minimum width=17mm,
                   minimum height=8mm, align=center, font=\scriptsize, line width=0.4pt},
  msg/.style    = {draw=black!60, fill=green!8, rounded corners=1.5pt, minimum width=15mm,
                   minimum height=7mm, align=center, font=\scriptsize, line width=0.4pt},
  flow/.style   = {-{Stealth[length=2mm]}, line width=0.7pt, draw=black!75},
  flowd/.style  = {-{Stealth[length=2mm]}, line width=0.7pt, draw=black!50, dashed},
  lab/.style    = {font=\scriptsize, inner sep=1.5pt, fill=white, fill opacity=0.85,
                   text opacity=1},
}
 载入，必须位于 \documentclass 之后。
%
% 编译：在 figures/ 目录下执行 xelatex fig1_architecture.tex
%      （含中文，必须用 xelatex；也可在 paper/ 下执行 make figures）
%
% 字体：默认使用 Windows 简体中文环境自带字体。
%      在 Linux/macOS 上请把下面三行替换为该系统已安装的 CJK 字体，例如
%      \setCJKmainfont{Noto Serif CJK SC}[BoldFont=Noto Sans CJK SC]
% =============================================================================

\usepackage{xeCJK}
\setCJKmainfont{SimSun}[BoldFont=SimHei]
\setCJKsansfont{Microsoft YaHei}

\usepackage{tikz}
\usepackage{amsmath,amssymb}
\usetikzlibrary{arrows.meta,positioning,fit,backgrounds,calc,shapes.geometric,decorations.pathreplacing}

% --- 统一样式 ---------------------------------------------------------------
\tikzset{
  dev/.style    = {draw=black!70, fill=blue!6,  rounded corners=2pt, minimum width=26mm,
                   minimum height=11mm, align=center, font=\small, line width=0.5pt},
  devr/.style   = {draw=black!70, fill=orange!10, rounded corners=2pt, minimum width=26mm,
                   minimum height=11mm, align=center, font=\small, line width=0.5pt},
  blk/.style    = {draw=black!60, fill=gray!8, rounded corners=1.5pt, minimum width=17mm,
                   minimum height=8mm, align=center, font=\scriptsize, line width=0.4pt},
  msg/.style    = {draw=black!60, fill=green!8, rounded corners=1.5pt, minimum width=15mm,
                   minimum height=7mm, align=center, font=\scriptsize, line width=0.4pt},
  flow/.style   = {-{Stealth[length=2mm]}, line width=0.7pt, draw=black!75},
  flowd/.style  = {-{Stealth[length=2mm]}, line width=0.7pt, draw=black!50, dashed},
  lab/.style    = {font=\scriptsize, inner sep=1.5pt, fill=white, fill opacity=0.85,
                   text opacity=1},
}


\begin{document}
\begin{tikzpicture}[x=1mm, y=1mm]

  % 泳道宽度与块高
  \def\W{150}
  \def\Hc{11}

  % 标题
  \node[font=\small, anchor=west] at (0, 30)
       {\textbf{异步流水：接收 $\parallel$ 计算 $\parallel$ 发送}};

  % --- 泳道 1：通信流（接收） ---
  \node[font=\scriptsize, anchor=east] at (-2, 20) {通信流};
  \draw[draw=black!50] (0,20) -- (\W,20);
  \foreach \i/\a in {0/0, 1/42, 2/84, 3/126}{
    \fill[blue!25, draw=blue!55, line width=0.3pt, rounded corners=1pt]
      (\a, 15.5) rectangle ++(38, \Hc);
    \node[font=\scriptsize] at (\a+19, 21) {recv $C^{s_{\i}}$};
  }
  % 发送流（与接收并行）
  \foreach \i/\a in {0/20, 1/62, 2/104}{
    \fill[orange!28, draw=orange!60, line width=0.3pt, rounded corners=1pt]
      (\a, 5) rectangle ++(38, \Hc);
    \node[font=\scriptsize] at (\a+19, 10.5) {send $C_{\i\to t}$};
  }

  % --- 泳道 2：计算流 ---
  \node[font=\scriptsize, anchor=east] at (-2, -14) {计算流};
  \draw[draw=black!50] (0,-14) -- (\W,-14);
  % 计算块按到达时间错开，宽 36、间隔 38 —— 形成连续流水
  \foreach \i/\a in {0/0, 1/38, 2/76, 3/114}{
    \fill[green!26, draw=green!60!black, line width=0.3pt, rounded corners=1pt]
      (\a, -20) rectangle ++(36, \Hc);
    \node[font=\scriptsize] at (\a+18, -9.5) {Attn $s_{\i}$};
  }

  % --- 到达事件箭头（块到达即刻触发计算） ---
  \foreach \a in {38, 76, 114}{
    \draw[-{Stealth[length=1.6mm]}, dashed, draw=red!65, line width=0.4pt]
      (\a, 15.5) -- (\a, -9.5);
  }
  \node[font=\scriptsize, text=red!65, anchor=west] at (116, 2)
       {到达即计算};

  % 时间轴
  \draw[-{Stealth[length=2mm]}, draw=black!60] (0,-26) -- (\W+6,-26)
       node[right, font=\scriptsize]{$t$};

  % --- 串行 vs 异步对比条 ---
  \begin{scope}[yshift=-40mm]
    \node[font=\scriptsize, anchor=east] at (-2, 4) {同步（串行）};
    \foreach \i/\a in {0/0, 1/30, 2/60, 3/90}{
      \fill[blue!18, draw=blue!40, line width=0.3pt] (\a, 0) rectangle ++(18, 7);
      \fill[green!18, draw=green!50!black, line width=0.3pt] (\a+18, 0) rectangle ++(18, 7);
    }
    \draw[draw=black!35] (0,0) rectangle (126, 7);
    \node[font=\scriptsize, anchor=west] at (130, 3.5) {$\approx 4\,(T_c+T_k)$};

    \node[font=\scriptsize, anchor=east] at (-2, -13) {异步（流水）};
    % 通信流整体：4 × T_c
    \fill[blue!18, draw=blue!40, line width=0.3pt] (0,-17) rectangle ++(120, 7);
    % 计算流：首个计算需等首个块到达(T_c)，之后每 T_k 一个，共 4 个
    \fill[green!30, draw=green!50!black, line width=0.5pt] (30,-17) rectangle ++(90, 7);
    \node[font=\scriptsize, anchor=west] at (126, -13.5)
         {$\approx \max(4T_c,\, T_c + 4T_k)$};
  \end{scope}

\end{tikzpicture}
\end{document}
